Addressing these limitations requires future work on dedicated long-closure detection, broader benchmarking, and integration with learned detectors. A dedicated long-closure module should first score normal blinks and prolonged closures separately. This distinction is necessary because the events most directly associated with drowsiness are sustained eye closures rather than ordinary blinks, and aggregating both categories under a single detection criterion may obscure the signal of greatest operational relevance~\citep{tian2021fatigue,shahbakhti2023fusion}. The benchmark should also be extended beyond driving data to other epoch-structured paradigms, including settings in which task blocks, trials, or recording segments provide natural units for estimating a representative signal distribution. Such evaluation would determine whether epoch screening remains useful when blink morphology, cognitive demands, and recording conditions differ from those in sustained driving. A further direction is to combine the screening stage with learned blink and drowsiness detectors rather than treating the approaches as alternatives~\citep{cabot2024drowsiness,makhmudov2024real,zhang2025fatigue}. In that configuration, epoch screening would serve as a computationally inexpensive front end that limits implausible candidate regions before a learned model performs the more specific classification.
